\section{Kode Antara: Postfix, N-Tuple, dan Three-Address Code}
\label{sec:ch14-kode-antara}

\subsection{Dari Aturan Semantik ke Kode Antara}

Aturan semantik tidak berhenti pada tipe. Setelah pohon sintaks
dianotasi, aturan yang sama dapat menghasilkan kode antara, yaitu
representasi program untuk mesin abstrak. Guntoro menempatkan kode
antara, notasi postfix, dan notasi n-tuple sebagai keluaran analisis
semantik; Juarna merinci bentuk three-address code. Ketiga bentuk itu
menjadi jembatan menuju pembangkitan kode target.

\subsection{Notasi Postfix}

Pada notasi postfix, operator ditulis setelah operand. Ekspresi
infix \((a+b)*(c+d)\) menjadi
\[
a\ b\ + \ c\ d\ + \ *.
\]
Evaluasi dilakukan dengan stack: operand didorong ke stack; ketika
operator ditemukan, dua operand teratas diambil, dioperasikan, lalu
hasilnya didorong kembali. Notasi ini tidak memerlukan tanda kurung
karena urutan operasi sudah ditentukan posisi operator.

\begin{table}[H]
\centering
\caption{Perbandingan infix dan postfix}
\label{tab:ch14-postfix}
\begin{tabularx}{\textwidth}{@{}l l X@{}}
\toprule
\textbf{Infix} & \textbf{Postfix} & \textbf{Catatan} \\
\midrule
\code{a + b} & \code{a b +} & satu operator \\
\code{a + b * c} & \code{a b c * +} & perkalian lebih dahulu \\
\code{(a + b) * c} & \code{a b + c *} & kurung mengubah urutan \\
\code{a = b + c * d} & \code{b c d * + a =} & assignment sebagai operator \\
\bottomrule
\end{tabularx}
\end{table}

\subsection{Notasi N-Tuple}

Notasi n-tuple menulis instruksi sebagai tuple. Dua bentuk yang umum
adalah \emph{triple} dan \emph{quadruple}. Triple memuat operator dan
dua operand; hasil antara dirujuk dengan nomor triple. Quadruple
menambahkan medan hasil secara eksplisit.

Untuk pernyataan \code{a = b + c * d}, quadruple-nya adalah:
\begin{center}
\begin{tabular}{@{}c l l l l@{}}
\toprule
\textbf{No} & \textbf{Operator} & \textbf{Operand 1} &
\textbf{Operand 2} & \textbf{Hasil} \\
\midrule
1 & \code{*} & \code{c} & \code{d} & \code{t1} \\
2 & \code{+} & \code{b} & \code{t1} & \code{t2} \\
3 & \code{=} & \code{t2} & --- & \code{a} \\
\bottomrule
\end{tabular}
\end{center}

Triple untuk pernyataan yang sama menulis hasil antara sebagai nomor:
\((1)\ *\ c\ d\), \((2)\ +\ b\ (1)\), dan
\((3)\ =\ (2)\ a\). Quadruple lebih mudah dibaca, sedangkan triple
lebih hemat ruang karena tidak menyimpan nama hasil antara.

\subsection{Three-Address Code}

Three-address code menulis setiap instruksi dengan paling banyak tiga
alamat: dua operand dan satu hasil. Bentuk ini dekat dengan kode
assembly dan mudah diterjemahkan ke instruksi mesin.
\[
\begin{aligned}
t_1 &= c * d\\
t_2 &= b + t_1\\
a &= t_2
\end{aligned}
\]
Setiap nama sementara hanya dipakai sekali. Urutan instruksi harus
menghormati dependensi: \(t_1\) dihitung sebelum \(t_2\), dan \(t_2\)
sebelum assignment ke \(a\).

\subsection{Kaitan dengan SDD dan Translation Scheme}

SDD dapat memuat atribut \code{code} yang menyimpan rangkaian
instruksi. Untuk produksi \(E \to E_1 + E_2\), aturan semantiknya
dapat berbentuk
\[
E.code = E_1.code \,\|\, E_2.code \,\|\,
\text{emit}(\code{+}, E_1.addr, E_2.addr, E.addr).
\]
Translation scheme menuliskan aksi yang sama pada posisi setelah
kedua operand selesai diproses. Dengan cara ini, tabel SDD pada
bagian sebelumnya dapat diuji tidak hanya pada tipe, tetapi juga pada
keluaran postfix atau three-address code.
